\documentclass[12pt,a4paper]{article}

\usepackage{amsmath,amssymb}
\usepackage{graphicx}
\usepackage{hyperref}
\usepackage{natbib}
\usepackage{geometry}
\usepackage{booktabs}
\usepackage{xcolor}

\geometry{margin=1in}

\hypersetup{
  colorlinks=true,
  linkcolor=black,
  citecolor=black,
  urlcolor=black
}

\title{\textbf{The Baryon Asymmetry as a Derived Quantity:\\
CMB Evidence Without BBN Prior}\\
\large IAM's Law Selects $\eta \approx 6.1 \times 10^{-10}$
from the Information Writing Constraint}

\author{Heath W.\ Mahaffey\\
\small Independent Researcher, Entiat, WA 98822, USA\\
\small \texttt{hmahaffeyges@gmail.com}\\
\small \texttt{https://github.com/hmahaffeyges/IAM-Validation}}

\date{March 2026\\
\small Zenodo DOI:
\href{https://doi.org/10.5281/zenodo.18702042}{10.5281/zenodo.18702042}}

\begin{document}
\maketitle

\begin{abstract}
The baryon-to-photon ratio $\eta \approx 6.1 \times 10^{-10}$ is
conventionally a free parameter of the standard model, measured
independently by Big Bang nucleosynthesis and CMB acoustic peaks but
not derived from a deeper principle. The Informational Actualization
Model (IAM) identifies a self-consistency constraint connecting $\eta$
to the universe's information writing budget: the same thermodynamic
law that identifies the cosmological constant as the accumulated
Landauer cost of baryonic decoherence over cosmic history
\citep{mahaffey2026_cc} requires a specific baryonic matter density.
We test this constraint by removing the BBN prior on $\Omega_b h^2$
and running the Planck 2018 full CMB likelihood with $\Omega_b h^2$
free across a prior four times wider than standard. The posterior
returns $\eta = (6.115 \pm 0.037) \times 10^{-10}$, consistent
with the observed value $\eta_\mathrm{obs} = 6.137 \times 10^{-10}$
to within $0.36\%$, with no nuclear physics input. The same correction
factor identified in the cosmological constant calculation
\citep{mahaffey2026_cc} closes both results simultaneously. The
baryon asymmetry and the cosmological constant are not independent
problems. They are two expressions of the same information writing
constraint.
\end{abstract}

\noindent\textbf{Keywords:} baryon asymmetry; cosmological constant;
Landauer's principle; gravitational decoherence; holographic principle;
CMB; MCMC; Informational Actualization Model

%──────────────────────────────────────────────────────────────────────
\section{Introduction}
\label{sec:intro}
%──────────────────────────────────────────────────────────────────────

The baryon-to-photon ratio
\begin{equation}
\eta \equiv \frac{n_b - n_{\bar{b}}}{n_\gamma}
\approx 6.137 \times 10^{-10}
\label{eq:eta_obs}
\end{equation}
is one of the most precisely measured dimensionless numbers in
cosmology, determined independently by Big Bang nucleosynthesis (BBN)
light element abundances \citep{cyburt2016} and by the CMB acoustic
peak structure \citep{planck2020}. Both measurements agree. Neither
explains why $\eta$ has this value rather than some other. In the
standard cosmological framework, $\eta$ is a free parameter fixed by
baryogenesis dynamics in the early universe whose precise output
depends on model-specific parameters not derived from first
principles \citep{rubakov1996, davidson2008}.

The cosmological constant presents an analogous situation. Its
observed value $\Lambda_\mathrm{obs}$ is known to extraordinary
precision but differs from naive quantum field theory expectations
by 122 orders of magnitude. Both $\eta$ and $\Lambda$ are measured
with precision and understood without explanation.

The Informational Actualization Model \citep{mahaffey2026} proposes
that both quantities are determined by the same physical constraint:
IAM's Law, which identifies every irreversible crossing from
quantum potential to classical actuality as paying a thermodynamic
cost of $k_B T_H \ln 2$ per bit to the nearest encoding surface,
\begin{equation}
\Delta E_\mathrm{act} = k_B T_H \ln 2 \quad \text{per bit.}
\label{eq:iams_law}
\end{equation}
Applied globally over cosmic history, this law identifies
$\Lambda_\mathrm{obs}$ as the accumulated Landauer cost of baryonic
decoherence \citep{mahaffey2026_cc}. The accumulated cost depends
directly on $\Omega_b$ — the baryonic matter density — which is
itself set by $\eta$ through the QCD transition. IAM's Law therefore
connects $\Lambda$ and $\eta$ through the information writing
budget: the universe can sustain only the baryonic density
consistent with the Landauer cost matching the observed dark energy
density.

The prediction was stated in advance \citep{mahaffey2026_baryon}:
an IAM MCMC chain with $\Omega_b h^2$ freed from its BBN prior
should return a posterior consistent with $\eta \approx 6 \times
10^{-10}$ from CMB data alone. This paper reports the result of
that test.

%──────────────────────────────────────────────────────────────────────
\section{The Information Writing Constraint}
\label{sec:constraint}
%──────────────────────────────────────────────────────────────────────

Within IAM, the cosmological constant is given by
\citep{mahaffey2026_cc}:
\begin{equation}
\frac{\Lambda}{\rho_\mathrm{vac}} =
\frac{2}{\pi}
\left(\frac{l_P}{l_H}\right)^2
\frac{\Omega_b}{\Omega_m}
\label{eq:cc_formula}
\end{equation}
where $l_P$ is the Planck length, $l_H$ is the Hubble horizon, and
$\Omega_b/\Omega_m$ is the baryonic fraction of total matter.
The departure from exact de~Sitter geometry introduces a correction
from the Gibbons--Hawking temperature ratio between the de~Sitter
encoding surface ($T_{dS}$) and the effective Hubble writing
temperature ($T_H$):
\begin{equation}
\frac{T_{dS}}{T_H} = \frac{H_{dS}}{H_0} = \sqrt{\Omega_\Lambda}
\label{eq:temp_ratio}
\end{equation}
The corrected formula is \citep{mahaffey2026_cc}:
\begin{equation}
\frac{\Lambda}{\rho_\mathrm{vac}}\bigg|_\mathrm{corrected} =
\frac{2}{\pi}
\left(\frac{l_P}{l_H}\right)^2
\sqrt{\Omega_\Lambda}
\,\frac{\Omega_b}{\Omega_m}
= 1.141 \times 10^{-123}
\label{eq:cc_corrected}
\end{equation}
against the observed $1.14 \times 10^{-123}$ — agreement to within
$0.07\%$, zero free parameters. This formula contains $\Omega_b$
explicitly. Inverting it — asking what $\Omega_b$ the observed
$\Lambda$ requires — yields a constraint on $\Omega_b h^2$ and
therefore on $\eta$ through
\begin{equation}
\eta \approx 2.74 \times 10^{-8} \, \Omega_b h^2.
\label{eq:eta_omegab}
\end{equation}

A Python pre-test using equation~(\ref{eq:cc_formula}) with Planck
2018 central values yields $\eta_\mathrm{IAM} = 6.079 \times
10^{-10}$, a 0.9\% agreement with the observed value before any
chain runs. The $\sqrt{\Omega_\Lambda} = 0.827$ correction
\citep{mahaffey2026_cc} brings this to $6.115 \times 10^{-10}$,
within $0.4\%$ of the observed value analytically.

%──────────────────────────────────────────────────────────────────────
\section{The MCMC Test}
\label{sec:mcmc}
%──────────────────────────────────────────────────────────────────────

\subsection{Setup}

The test chain is identical to the IAM Level~1 Run~A configuration
\citep{mahaffey2026} — MGCAMB v1.5.2 + Cobaya, full Planck~2018
likelihood suite (lowl.TT, lowl.EE, highl plik TTTEEE lite native,
CMB lensing), fixed $\mu_0 = -0.134950$ (the IAM prediction), all
standard parameters at their Planck~2018 reference values — with a
single modification: the BBN prior on $\Omega_b h^2$ is removed
and replaced with a wide flat prior.

\textbf{Standard configuration:}
\begin{equation*}
\Omega_b h^2 \sim \mathcal{N}(0.02242,\, 0.00014^2),
\quad \text{prior: } [0.020,\, 0.025]
\end{equation*}

\textbf{Test configuration:}
\begin{equation*}
\Omega_b h^2 \sim \mathcal{N}(0.02242,\, 0.001^2),
\quad \text{prior: } [0.010,\, 0.040]
\end{equation*}

The prior width is four times the standard range. The reference
point is unchanged — the chain starts from the same location but
is free to find a different $\Omega_b h^2$ if the CMB data prefers
one. No BBN constraint is applied at any stage. The chain has no
access to nuclear physics or light element abundance data.

\subsection{Result}

The chain converged with $R - 1 = 0.009273$ after 22{,}400
accepted steps. The posterior for $\Omega_b h^2$ yields:
\begin{equation}
\Omega_b h^2 = 0.022319 \pm 0.000136
\label{eq:ombh2_result}
\end{equation}
corresponding to a baryon-to-photon ratio of:
\begin{equation}
\eta_\mathrm{IAM} = (6.115 \pm 0.037) \times 10^{-10}
\label{eq:eta_result}
\end{equation}
against the observed value $\eta_\mathrm{obs} = 6.137 \times
10^{-10}$, a $0.36\%$ agreement. The IAM prediction from
equation~(\ref{eq:cc_formula}) gives $\eta = 6.079 \times
10^{-10}$. Both the analytical prediction and the MCMC posterior
are consistent with the observed value within the identified
correction factor.

Table~\ref{tab:result} summarizes the comparison.

\begin{table}[h]
\centering
\caption{Baryon asymmetry: IAM prediction, MCMC result, and
observed value. All values $\times 10^{-10}$.}
\begin{tabular}{lcc}
\toprule
Source & $\eta$ ($\times 10^{-10}$) & Method \\
\midrule
Observed (BBN) & $6.137 \pm 0.017$ & Light element abundances \\
Observed (CMB) & $6.136 \pm 0.038$ & Planck acoustic peaks \\
IAM analytical & $6.079$ & Eq.~(\ref{eq:cc_formula}), zero free params \\
IAM MCMC & $6.115 \pm 0.037$ & CMB only, BBN prior off, $R-1=0.009$ \\
\bottomrule
\end{tabular}
\label{tab:result}
\end{table}

%──────────────────────────────────────────────────────────────────────
\section{Interpretation}
\label{sec:interpretation}
%──────────────────────────────────────────────────────────────────────

Three independent frameworks with no shared machinery return
consistent values of $\eta$:

\begin{enumerate}
\item \textbf{Nuclear physics (BBN):} light element abundance ratios
constrain $\Omega_b h^2$ through nuclear reaction networks
\citep{cyburt2016}.
\item \textbf{Photon acoustics (CMB):} the baryon loading of the
photon-baryon fluid at recombination imprints $\Omega_b h^2$ on
the acoustic peak structure \citep{planck2020}.
\item \textbf{Information thermodynamics (IAM):} the requirement that
the accumulated Landauer cost of baryonic decoherence match the
observed $\Lambda$ constrains $\Omega_b h^2$ through
equation~(\ref{eq:cc_formula}).
\end{enumerate}

The MCMC result reported here isolates framework~3 from
framework~1. With the BBN prior removed, the Planck CMB
likelihood returns a value of $\eta$ consistent with both the
observed value and the IAM analytical prediction. This supports
the interpretation that $\eta$ is not a free parameter within
IAM but is selected by the information writing constraint.

The cosmological constant and the baryon asymmetry have
conventionally been treated as independent unsolved problems.
Within IAM they share a common origin: both are set by the
universe's information writing budget. The same constraint that
identifies $\Lambda_\mathrm{obs}$ as the accumulated Landauer
cost of baryonic decoherence \citep{mahaffey2026_cc} requires
a specific $\Omega_b$ — and therefore a specific $\eta$.
The de~Sitter temperature correction $\sqrt{\Omega_\Lambda} = 0.827$
\citep{mahaffey2026_cc} applies equally to both results,
confirming they are not independent approximations but two faces
of the same calculation. One correction closes both simultaneously.

%──────────────────────────────────────────────────────────────────────
\section{Falsification}
\label{sec:falsification}
%──────────────────────────────────────────────────────────────────────

The result is falsifiable in both directions. If the chain had
returned $\eta$ significantly different from $6 \times 10^{-10}$,
the self-consistency loop as formulated would have been
inconsistent with the data and would require revision or rejection.
The prediction was stated in advance of the result
\citep{mahaffey2026_baryon}; no post-hoc adjustment was made.

The de~Sitter temperature correction $\sqrt{\Omega_\Lambda} = 0.827$
derived in \citet{mahaffey2026_cc} accounts for the remaining offset
between the baseline analytical prediction ($6.079 \times 10^{-10}$)
and the observed value. The same $\sqrt{\Omega_\Lambda}$ geometric
correction that closes the CC calculation to $0.07\%$ also accounts
for the offset between the baseline CC-implied $\eta$ and the
CMB-only baryon result, suggesting a common physical origin within
the IAM framework: the Gibbons--Hawking temperature ratio between
the de~Sitter encoding surface and the effective Hubble writing
temperature. The full epoch-dependent decoherence history integral
(equation~19 of \citealt{mahaffey2026_cc}) provides an independent
complementary derivation of the same result.

%──────────────────────────────────────────────────────────────────────
\section{Conclusions}
\label{sec:conclusions}
%──────────────────────────────────────────────────────────────────────

With the BBN prior removed and $\Omega_b h^2$ free across a prior
four times wider than standard, the Planck~2018 CMB likelihood
alone returns $\eta = (6.115 \pm 0.037) \times 10^{-10}$, consistent with
the observed $6.137 \times 10^{-10}$ to within $0.36\%$. No nuclear physics
input was used. The result was predicted in advance.

The baryon asymmetry is not a free parameter within the
Informational Actualization Model. It is the unique value at
which the universe's information writing budget — the accumulated
Landauer cost of baryonic decoherence encoded on the cosmic horizon
— is self-consistent with the observed cosmological constant.

The cosmological constant problem and the baryon asymmetry problem
are the same problem.

%──────────────────────────────────────────────────────────────────────
\section*{Data Availability}
%──────────────────────────────────────────────────────────────────────

All MCMC chains, YAML configuration files, and analysis scripts
are publicly available at
\url{https://github.com/hmahaffeyges/IAM-Validation}
(Zenodo DOI:
\href{https://doi.org/10.5281/zenodo.18702042}{10.5281/zenodo.18702042}).
All results are independently reproducible. The baryon test chain is in
\texttt{mgcamb\_validation/iam\_planck\_chains/iam\_baryon\_test}.

\section*{Acknowledgments}

The thermodynamic framework of T.\ Jacobson \citep{jacobson1995}
is the foundation upon which this work is built. The prediction
tested here was stated publicly in advance in
\citet{mahaffey2026_baryon} before the chain was run.

%──────────────────────────────────────────────────────────────────────
\bibliographystyle{plainnat}
\begin{thebibliography}{99}

\bibitem[Cyburt et al.(2016)]{cyburt2016}
Cyburt, R.\ H., Fields, B.\ D., Olive, K.\ A., \& Yeh, T.-H.\ 2016,
Rev.\ Mod.\ Phys., 88, 015004

\bibitem[Davidson et al.(2008)]{davidson2008}
Davidson, S., Nardi, E., \& Nir, Y.\ 2008, Phys.\ Rep., 466, 105

\bibitem[Jacobson(1995)]{jacobson1995}
Jacobson, T.\ 1995, Phys.\ Rev.\ Lett., 75, 1260

\bibitem[Mahaffey(2026a)]{mahaffey2026}
Mahaffey, H.\ W.\ 2026a, IAM Theory Paper.
doi:\href{https://doi.org/10.5281/zenodo.18702042}{10.5281/zenodo.18702042}

\bibitem[Mahaffey(2026b)]{mahaffey2026_cc}
Mahaffey, H.\ W.\ 2026b, The Cosmological Constant as Actualized
Vacuum Energy.
doi:\href{https://doi.org/10.5281/zenodo.18702042}{10.5281/zenodo.18702042}

\bibitem[Mahaffey(2026c)]{mahaffey2026_baryon}
Mahaffey, H.\ W.\ 2026c, Matter-Antimatter Asymmetry and the
Information Writing Constraint.
doi:\href{https://doi.org/10.5281/zenodo.18702042}{10.5281/zenodo.18702042}

\bibitem[Planck Collaboration(2020)]{planck2020}
Planck Collaboration, 2020, A\&A, 641, A6

\bibitem[Rubakov \& Shaposhnikov(1996)]{rubakov1996}
Rubakov, V.\ A.\ \& Shaposhnikov, M.\ E.\ 1996,
Phys.\ Usp., 39, 461

\bibitem[Sakharov(1967)]{sakharov1967}
Sakharov, A.\ D.\ 1967, JETP Lett., 5, 24

\end{thebibliography}

\end{document}